\subsection{HD-FT03: split alternating constituents as arbitrary tensor factors}
\label{sec:split-alternating-factor}

\paragraph{Status and provenance.}
This is a new addendum by \texttt{/root/imprimitive} to the root
HD-FT01--02 argument.  The symmetric-corner construction and the
explicit inequalities below are new derivations.  The alternating
character formula is a published input, with exact locators below.
Complete author and separate root whole-proof readings pass. Root
also read the complete source/status record, bibliography, author
checker and first successful output; the exact reading and source
access boundaries are recorded in the retained root review.
No projected twelve-mark hypothesis is introduced.

\begin{theorem}[HD-FT03]
Let \(\lambda=\lambda'\vdash n\), and let \(W\) be either of the
two ordinary irreducible constituents of the restriction of the
\(S_n\)-module labelled by \(\lambda\) to \(A_n\).
Write \(a=\dim W\), and let \(f(x)\) count the fixed natural letters
of \(x\in A_n\).  If \(n\ge65536\), then
\[
 a\ge n^{120},\qquad
 \frac{|\chi_W(x)|}{a}\le
 \begin{cases}
       (f(x)/n)^3,&f(x)>0,\\
       n^{-6},&f(x)=0 .
 \end{cases}
 \tag{FTS.1}
\]
Consequently, when \(n>3500000\), such a \(W\) cannot occur in the
actual factorwise tensor configuration (FT.1)--(FT.2) of HD-FT01,
with arbitrary tensor partner and arbitrary scalar factors.
This conclusion does not depend on the original imaginary
quadratic field or on the number of nonidentity projected inertias.
\end{theorem}

\begin{proof}
Put \(F=f^\lambda\), the degree of the \(S_n\)-module, so that
\(a=F/2\).  Let \(r\) be the number of diagonal boxes of \(\lambda\),
and write its diagonal hook lengths as
\[
 d_1>d_2>\cdots>d_r>0,\qquad
 d_i=2(\lambda_i-i)+1,\qquad
 \sum_i d_i=n,\qquad r^2\le n.
 \tag{FTS.2}
\]
The hooks in (FTS.2) are distinct odd positive integers.  The sum
identity follows by partitioning the symmetric diagram into its
diagonal hooks, and their minimum possible sum is \(r^2\).

First we give a lower bound for \(F\) using standard tableaux.
Starting with any nonempty self-conjugate diagram, select a corner
by a fixed deterministic rule.  If it is off the diagonal, its
transpose is a distinct corner.  These two boxes are incomparable
in the diagram order.  Delete both, assigning to them the two
largest available labels in either order.  If the selected corner
is diagonal, delete it and assign the largest available label.
After either operation the remaining diagram is self-conjugate.
The two choices in a paired deletion do not change this remaining
diagram, and therefore do not change any subsequent selected pair.
Reversing the deletions produces standard tableaux: each deletion
is a valid corner removal, including both orders of a paired removal.
Different choices produce different labelled tableaux.
Exactly \(r\) diagonal boxes are removed singly, and all other
boxes are removed in pairs.  The usual interpretation of \(F\) as
the number of standard tableaux, equivalently repeated branching,
therefore gives
\[
                    F\ge 2^{(n-r)/2}
                       \ge2^{(n-\sqrt n)/2}.
                    \tag{FTS.3}
\]
This construction uses a subset of all tableaux and does not assert
an exact formula for their number.

We record elementary logarithmic estimates with a completely
explicit endpoint.  For every real \(x\ge2^{16}\),
\[
 1+\log_2x\le\frac{\sqrt x}{4},\qquad
 1+6\log_2x\le\frac{x}{4}.
 \tag{FTS.4}
\]
At \(x=2^{16}\) these read \(17\le64\) and \(97\le16384\).
The differences between the right and left sides have derivatives
\[
 \frac1{8\sqrt x}-\frac1{x\log2}
   \ge\frac1{8\sqrt x}-\frac2x\ge0,\qquad
 \frac14-\frac6{x\log2}\ge\frac14-\frac{12}x\ge0.
 \tag{FTS.5}
\]
Here \(\log2=\int_1^2dt/t\ge1/2\), and the displayed ranges imply
the final inequalities.  This proves (FTS.4) on the full half-line,
not merely at integer or dyadic sample points.
In particular, writing \(s=\sqrt n\ge256\), (FTS.3)--(FTS.4) give
\[
 \log_2(F/2)-120\log_2n
 \ge\frac{n-s}{2}-1-(30s-120)
 =\frac{s(s-61)}2+119>0.
 \tag{FTS.6}
\]
Thus \(a=F/2\ge n^{120}\).

We next use the exact published alternating split-character formula.
If the cycle type of \(x\) differs from \((d_1,\ldots,d_r)\), then
\[
                 \chi_W(x)=\frac12\chi^\lambda(x).
                 \tag{FTS.7}
\]
At either of the two exceptional alternating classes of that
cycle type,
\[
 \chi_W(x)=\frac12\left(\varepsilon
                \mathbin{\pm}\sqrt{\varepsilon\prod_i d_i}\right),
 \qquad \varepsilon=(-1)^{(n-r)/2}.
 \tag{FTS.8}
\]
The choice of sign depends on the class and constituent and will
not be needed.  Exact primary locators are Balcerzyk,
\cite{FTS:Balcerzyk1997}, formulas (3)--(4), pp.~148--149, and the
diagonal-hook character value on p.~151.  The same restriction and
split-value formulas are restated in the published research paper
\cite{FTS:YangStaroletov2021}, Propositions 3--4, author pp.~5--6,
with the branching rule in Proposition 5.
One may also obtain the sign in (FTS.8) by Murnaghan--Nakayama:
remove the unique rim hook of largest diagonal-hook length at
each stage; the remaining inner diagram is self-conjugate, and
the total leg length is \(\sum_i(d_i-1)/2=(n-r)/2\).

For an exceptional element, (FTS.2)--(FTS.3) yield
\[
 \frac{|\chi_W(x)|}{a}
 \le\frac{1+\sqrt{\prod_i d_i}}{F}
 \le 2^{\,1-n/2+\sqrt n(1+\log_2n)/2}.
 \tag{FTS.9}
\]
Indeed \(\prod_i d_i\le n^r\) and \(r\le\sqrt n\).
Using the two estimates in (FTS.4) in sequence, the exponent in
(FTS.9) is at most
\[
 1-\frac{3n}{8}
 =\left(1-\frac n4\right)-\frac n8
 \le-6\log_2n-\frac n8
 <-6\log_2n .
 \tag{FTS.10}
\]
Thus every exceptional value has absolute ratio strictly less
than \(n^{-6}\).  Moreover its cycle lengths are distinct odd
integers, so it has at most one fixed letter.  If \(f(x)=1\), then
\(n^{-6}\le n^{-3}=(f(x)/n)^3\); if \(f(x)=0\), the required
fixed-point-free estimate holds directly.

For every nonexceptional \(x\), (FTS.7) shows that
\(|\chi_W(x)|/a=|\chi^\lambda(x)|/F\).
The already checked effective MSP consequence used in HD-TW02
and HD-FT02 applies because \(n\ge65536>4000\) and
\(F\ge2n^{120}\).  It gives the cubic bound for positive fixed-point
count and the stronger fixed-point-free bound
\[
                  |\chi^\lambda(x)|/F<F^{-1/20}
                                      \le n^{-6}.
                  \tag{FTS.11}
\]
The identity ratio is exactly one.  This proves all of (FTS.1).

Finally the bound (FTS.1) satisfies the weaker error allowance
\(2n^{-6}\) in (FT.3).  All hypotheses of HD-FT01 now hold for the
actual quotient \(A_n\) and actual tensor factorization, so that
the source is impossible for \(n>3500000\).
This last step uses actual quotient curves and their genera, not
a claim that the projected factor has twelve branch points.
\end{proof}

\paragraph{Compatible-normalizer precision.}
For \(n\ge7\), the known automorphism description of \(A_n\) identifies
its outer automorphism with conjugation by an odd permutation.
That automorphism exchanges the two inequivalent split constituents
(the same Proposition 3(ii) cited above).
It cannot stabilize the projective class of either one:
if an intertwiner implemented it up to scalar, the scalar discrepancy
would be a linear character of the perfect group \(A_n\), hence trivial.
Schur's lemma then shows that the compatible linear normalizer is
scalar times the represented \(A_n\).  This uses the previously checked
ordinary alternating automorphism input; it asserts no new
classification theorem.  Thus the \(A_n\)-quotient interface above
also covers a factor with its full compatible scalar normalizer.
As always, an actual quotient and tensor factorization must be supplied;
they are not inferred from a numerical spectrum alone.

\paragraph{Combined range.}
HD-FT02 and HD-FT03 together treat every ordinary irreducible
alternating constituent at degree at least \(n^{120}\).
In particular all split constituents disappear already for
\(n>3500000\), and for
\(n\ge241^{482}\) every surviving ordinary symmetric or alternating
factor has a nonsplit partition label with row or column tail at
most \(240\).  Arbitrary factorwise tensor partners are still allowed.
This does not settle these finitely many tail shapes in unbounded
natural degree, and does not assert a finite catalogue of all
projective factor types.
